\chapter{Practical System Design}

\section{Data Pipelines}

In modern machine learning systems, data is as important as models. Effective system design begins with reliable and scalable data pipelines.

\medskip

\noindent
\textbf{Data pipeline.}  
A data pipeline transforms raw data into a form suitable for training and evaluation.

\medskip

\noindent
\textbf{Stages.}
\begin{itemize}
    \item Data collection
    \item Cleaning and preprocessing
    \item Feature extraction or encoding
    \item Storage and access
\end{itemize}

\medskip

\noindent
\textbf{Preprocessing.}
\begin{itemize}
    \item Normalization and scaling
    \item Handling missing values
    \item Data augmentation
\end{itemize}

\medskip

\noindent
\textbf{Streaming vs batch processing.}
\begin{itemize}
    \item Batch pipelines process data in large chunks
    \item Streaming pipelines process data continuously
\end{itemize}

\medskip

\noindent
\textbf{Data quality.}
\begin{itemize}
    \item Consistency and correctness
    \item Noise and bias considerations
\end{itemize}

\medskip

\noindent
\textbf{Insight.}  
The quality of a machine learning system is fundamentally limited by the quality of its data.

\section{Training Workflows}

Training modern models requires coordinated workflows that manage data, computation, and experimentation.

\medskip

\noindent
\textbf{Training loop.}
\begin{enumerate}
    \item Load batch of data
    \item Compute predictions
    \item Evaluate loss
    \item Compute gradients
    \item Update parameters
\end{enumerate}

\medskip

\noindent
\textbf{Experiment management.}
\begin{itemize}
    \item Track hyperparameters
    \item Record metrics and logs
    \item Compare different model configurations
\end{itemize}

\medskip

\noindent
\textbf{Hyperparameter tuning.}
\begin{itemize}
    \item Grid search, random search
    \item Adaptive methods (e.g., Bayesian optimization)
\end{itemize}

\medskip

\noindent
\textbf{Distributed training.}
\begin{itemize}
    \item Data parallelism across multiple devices
    \item Model parallelism for large models
\end{itemize}

\medskip

\noindent
\textbf{Checkpointing.}
\begin{itemize}
    \item Save intermediate model states
    \item Enable recovery and experimentation
\end{itemize}

\medskip

\noindent
\textbf{Insight.}  
Well-designed workflows enable efficient experimentation and reliable model development.

\section{Debugging and Diagnostics}

Training machine learning models often involves diagnosing issues that arise during optimization.

\medskip

\noindent
\textbf{Common problems.}
\begin{itemize}
    \item Training does not converge
    \item Overfitting or underfitting
    \item Numerical instability
\end{itemize}

\medskip

\noindent
\textbf{Diagnostic tools.}
\begin{itemize}
    \item Training and validation curves
    \item Gradient norms
    \item Activation statistics
\end{itemize}

\medskip

\noindent
\textbf{Sanity checks.}
\begin{itemize}
    \item Verify loss decreases on small datasets
    \item Check gradient correctness
    \item Ensure labels and data are consistent
\end{itemize}

\medskip

\noindent
\textbf{Visualization.}
\begin{itemize}
    \item Monitor metrics over time
    \item Inspect intermediate outputs
\end{itemize}

\medskip

\noindent
\textbf{Systematic debugging.}
\begin{itemize}
    \item Isolate components
    \item Simplify the model
    \item Incrementally increase complexity
\end{itemize}

\medskip

\noindent
\textbf{Insight.}  
Debugging is an essential skill in machine learning, requiring both theoretical understanding and practical experience.

\section{Reproducibility}

Reproducibility is critical for scientific and engineering reliability.

\medskip

\noindent
\textbf{Definition.}  
A result is reproducible if it can be consistently obtained under the same conditions.

\medskip

\noindent
\textbf{Sources of variability.}
\begin{itemize}
    \item Random initialization
    \item Data shuffling
    \item Hardware and numerical differences
\end{itemize}

\medskip

\noindent
\textbf{Best practices.}
\begin{itemize}
    \item Fix random seeds
    \item Record configurations and hyperparameters
    \item Version datasets and code
\end{itemize}

\medskip

\noindent
\textbf{Experiment tracking.}
\begin{itemize}
    \item Store results systematically
    \item Maintain logs of experiments
\end{itemize}

\medskip

\noindent
\textbf{Determinism vs performance.}
\begin{itemize}
    \item Fully deterministic execution may reduce performance
    \item Tradeoff between reproducibility and efficiency
\end{itemize}

\medskip

\noindent
\textbf{Insight.}  
Reproducibility ensures that results are reliable and can be built upon.

\section{A Unifying Perspective}

Practical system design integrates multiple components:

\begin{itemize}
    \item \textbf{Data:} pipelines and preprocessing
    \item \textbf{Computation:} training workflows
    \item \textbf{Diagnostics:} debugging tools
    \item \textbf{Reliability:} reproducibility practices
\end{itemize}

\medskip

\noindent
\textbf{Core idea.}  
Successful machine learning systems require careful coordination of data, models, and infrastructure.

\medskip

\noindent
\textbf{Key insight.}  
Practical effectiveness depends not only on algorithms, but on the systems that support them.

\section{Outlook}

This chapter examined practical aspects of machine learning systems:
\begin{itemize}
    \item data pipelines,
    \item training workflows,
    \item debugging and diagnostics,
    \item reproducibility.
\end{itemize}

\medskip

This concludes the book.

\medskip

\noindent
\textbf{Final guiding principle.}  
Machine learning is both a mathematical discipline and an engineering practice, requiring insight into theory, data, and systems.